%!TEX root = ../../clcxsj.tex

\chapter{ \cs 基础语法}

\cs 程序设计语言是类 C 语言，基本语法与 C 或 \cpp 基本类似，
比如块注释为 \verb|/* */|， 行注释为 \verb|//| ，
每行语句用分号 \verb|;| 分隔，\cs 标识符区分大小写。
\cs 使用相同的大括号 \verb|{ }| 、 中括号 \verb|[ ]| 与 圆括号 \verb|( )|,
控件语句 \verb|if、else、switch|，
循环语句 \verb|for、while、do while|等。

本章主要讲解 \cs 中不同于  C 或 \cpp 的语法。


\section{ \cs 中的变量与数据类型}

为保持 \cs 语言的高效性，\cs 中的数据类型分为两类：值类型 (value type) 和引用类型 (reference type)。
从概念上讲，值类型变量直接存储其值，引用类型变量存储对其值的引用。

值类型变量的值通常直接在栈上保存，变量直接存储相应的值，因此变量的值自动生成、自动释放。
当值类型的变量进行赋值或传参操作时，会创建该值的副本，对一个值类型变量的修改不会影响另一个。
基本数据类型如 int、float、double、decimal、bool、char 与 结构体（struct）、枚举（enum）等都是值类型。

引用类型变量在堆中生成变量的值，在栈中保存堆中值的地址，当变量被释放时，
在堆中生成的值不会自动释放，靠 .NET 架构回收。
引用类型包括类（class）、接口（interface）、委托(delegate)、数组（array）以及字符串（string）等。

值类型和引用类型在内存中的存储方式差异导致了它们在 \cs 程序中的不同行为。
例如，值类型在赋值时会复制整个数据，而引用类型在赋值时只复制引用。
因此对引用类型的一个实例的修改会影响所有指向该实例的变量，而值类型的修改则限定于单个实例。

值类型适用于表示简单的数据结构，如数值和布尔值。值类型通常具有更好的性能，因为它们可以直接存储在栈上，
栈上的内存分配和回收速度非常快，在内存管理方面更高效，且不需要垃圾回收。
引用类型适用于表示更复杂的数据结构和对象，如字符串和数组等。它们支持继承和多态，但需要垃圾回收来管理内存。

\subsection{值类型}
\cs  的值类型进一步划分为基本类型、枚举 (enum)、结构(struct)
和可以为 null 的值类型 (nullable type)。

对于值类型，每个变量都有它们自己的数据副本（除 ref 和 out 参数变量外），
因此对一个变量的操作不可能影响另一个变量。

\begin{tabular}{l|l}
    \hline
    类别         & 说明                                   \\
    \hline
    基本类型       & 有符号整型：sbyte 、short 、int 、long        \\
               & 无符号整型：byte、ushort、uint、ulong         \\
               & Unicode 字符：char                      \\
               & IEEE 浮点：float、double                 \\
               & 高精度小数型：decimal                       \\
               & 布尔：bool                              \\
    结构(struct) & struct S {...} 形式的用户定义的类型            \\
    枚举(enum)   & enum E {...} 形式的用户定义的类型              \\
    可空(T?)     & 派生自 System.Nullable<T> 泛型结构体，T? 为其别名 \\
    \hline
\end{tabular}


\subsection{引用类型}
\cs  的引用类型进一步划分为类 (class)、接口(interface)、
数组 (array) 和委托(delegate)。
对于引用类型，两个变量可能引用同一个对象，因此对一个变量的操作可能影响另一个变量所引用的对象。

\begin{tabular}{l|l}
    \hline
    类别            & 说明                                                  \\
    \hline
    类(class)      & 派生于 System.Object, 其别名为 object ，是所有其他类型的基类          \\

    接口(interface) & interface I {...} 形式的用户定义的类型                        \\
    数组(Array)     & 派生于 System.Array, 一维和多维数组，例如 int[] 和 int[,]         \\
    委托(delegate)  & 派生于 System.Delegate, delegate int D(...) 形式的用户定义的类型 \\
    \hline
\end{tabular}

\cs 中的 string 的类名为 System.String，string 是其别名，因此 string 是引用类型数据。

\cs 中常用的 DateTime 的定义为 \verb|public readonly struct DateTime : ......|，因此DateTime是值类型数据。

\cs 的所有类型都继承自 System.Object 类，System.Object 是所有类型的基类，


\subsection{自定义类型}
\cs  程序使用类型声明 (type declaration) 创建新类型。类型声明指定新类型的名称和成员。
在  \cs  类型分类中，有五类是用户可定义的：类(class)、结构(struct)、
接口(interface)、枚举(enum) 和委托(delegate)。

注：\cs 7.0 提供了元组(tuple), \cs 9.0 提供了记录(record)。


类是定义了一个包含数据成员（字段）和函数成员（方法、属性等）的数据结构，类支持单一继承和多态。

结构与类相似，表示一个带有数据成员和函数成员的结构。
但是，与类不同，结构是值类型，不需要堆分配。
结构类型不支持用户指定的继承，所有结构类型都从类 System.ValueType 继承。

接口定义了一个协定，是一个公共函数成员的命名集。
实现某个接口的类或结构必须提供该接口的函数成员的实现。
一个接口可以从多个接口继承，而一个类或结构可以实现多个接口。

委托表示对具有特定参数列表和返回类型的方法的引用。
通过委托，我们能够将方法作为实体赋值给变量和作为参数传递。
委托类似于在其他某些语言中的函数指针的概念，但是与函数指针不同，委托是面向对象的，并且是类型安全的。

类、结构、接口和委托都支持泛型，因此可以通过其他类型将其参数化。

枚举是具有命名常量的独特的类型。每种枚举类型都具有一个基础类型，
该基础类型必须是八种整型之一。枚举类型的值集和它的基础类型的值集相同。

\cs  支持由任何类型组成的一维和多维数组。与以上列出的类型不同，数组类型不必声明就可以使用。
实际上，数组类型是通过在某个类型名后加一对方括号来构造的。
例如，int[] 是一维 int 数组，int[,] 是二维 int 数组，int[][] 是一维 int 数组的一维数组。

可以为 null 的类型也不必声明就可以使用。对于每个不可以为 null 的值类型 T，
都有一个相应的可以为 null 的类型 T?，该类型可以容纳附加值 null。
例如，int? 类型可以容纳任何 32 位整数或 null 值。

\subsection{装箱与拆箱}
\cs  的类型系统是统一的，因此任何类型的值都可以按对象处理。 \cs  中的每个类型直接或间接地从 object 类类型派生，
而 object 是所有类型的根基类。引用类型的值都被视为 object 类型，被简单地当作对象来处理。
值类型的值则通过对其执行装箱和拆箱操作按对象处理。下面的示例将 int 值转换为 object，然后又转换回 int。

\begin{lstlisting}
using System;
class Test
{
    static void Main()
    {
        int i = 123;
        object o = i;      // Boxing
        int j = (int)o;    // Unboxing
    }
}
\end{lstlisting}

当将值类型的值转换为类 object 时，将分配一个对象实例（也称为“箱子”）以包含该值，
并将值复制到该箱子中。反过来，当将一个 object 引用强制转换为值类型时，
将检查所引用的对象是否含有正确的值类型，如果有，则将箱子中的值复制出来。

在程序中频繁地装箱和拆箱操作会带来性能开销，应尽量避免不必要的装箱和拆箱操作。

\subsection{延伸阅读值类型}
所有值类型隐式继承自 System.ValueType ， 而 System.ValueType 又继承自 System.Object 类。
这听起来是不是有点自我矛盾？其中有什么奥秘？

实际上值类型和引用类型都是从 System.Object 继承的，System.Object 类是 \cs 中所有类的基类。
所有值类型都隐式继承自类 System.ValueType，而类 System.ValueType 又继承自类 System.Object（别名object），
奥秘的地方在于 System.ValueType 用更合适的值类型实现重写了 System.Object 中的虚方法。

但任何类型都不能直接从值类型 System.ValueType 派生。

这样就有一个问题，值类型与引用类型都是从 System.Object 继承的，怎么 object 类型与 int 类型就不一样呢？

这里大致可以这样理解，.NET 是动态加载运行 \cs 代码的，CLR 在运行时把值类型装载到栈，把引用类型装载到堆，
栈中只提供一个引用指向堆，即引用。

有关具体的值类型与引用类型的区别，请参考有关微软官方文档。

\subsection{\cs 中较为特别的类型}

\begin{itemize}
    \item 元组（tuple）

          元组是 \cs 中一种轻量级的数据结构，用于将多个值组合成一个复合值。
          它在 \cs 4.0 中首次引入，随后在 \cs 7.0 中通过 ValueTuple 提供了更强大的功能。
          元组可以存储不同类型的元素，并且无需显式定义字段类型和名称。

          如测绘中常用的六十进制的角度值可以用显示的元组表示为：

          \begin{lstlisting}
(int d, int m, double s) a1 = (45, 30, 15.5);
Console.WriteLine($"Angle value is {a1.d}°{a1.m}'{a1.s}\".");
// Output: Angle value is 45°30'15.5".
\end{lstlisting}

          也可以隐式的定义为：

          \begin{lstlisting}
var a1 = (45, 30, 15.5);
Console.WriteLine($"Angle value is {a1.Item1}°{a1.Item2}'{a1.Item3}\".");
// Output: Angle value is 45°30'15.5".
\end{lstlisting}

          元组类型是值类型，元组元素是公共字段, 若无显示类型指定， 元组字段的默认名称为 Item1、Item2、Item3 等。

          还可以定义为：

          \begin{lstlisting}
var (d, m, s) = (45, 30, 15.5);
Console.WriteLine($"Angle value is {d}°{m}'{s}\".");
// Output: Angle value is 45°30'15.5".
\end{lstlisting}

          上面代码利用 var 关键字，由编译器根据右边的值来推断变量 d、m、s 的类型， 从而保持代码的简洁。



    \item 记录（record）

          record 修饰符是用于封装数据的内置功能，提供了一种简洁的语法来定义不可变的数据对象。
          record 与 record class 语法定义引用类型，record struct 定义值类型。

          如下代码定义了一个 record Point：
          \begin{lstlisting}
public record Point(string Name, double X, double Y, double Z);
\end{lstlisting}

          在 record 上声明主构造函数时，编译器会为主构造函数参数生成公共属性。
          上面代码相当于如下代码：
          \begin{lstlisting}
public record Point
{
    public required string Name { get; init; }
    public required double X { get; init; }
    public required double Y { get; init; }
    public required double Z { get; init; }
};
\end{lstlisting}

          其中 init 修饰符表示属性只能在对象初始化时设置值，之后不能修改。

          record 类型的对象是不可变的，意味着一旦创建了一个 record 对象，其属性值就不能被修改。
          这有助于确保数据的一致性和安全性，尤其是在多线程环境中。

          我们在Top-Level Statement中测试使用record类型，如下代码所示：
          \begin{lstlisting}
Point point1 = new("GP10", 3804770.542, 608960.290, 471.823);
// point1.Name = "GP11";
Console.WriteLine(point1);
// output: Point { Name = GP10, X = 3804770.542, Y = 608960.290, Z = 471.823 }

Point point2 = new("GP10", 3804770.542, 608960.290, 471.823);
Console.WriteLine(point2);
// output: Point { Name = GP10, X = 3804770.542, Y = 608960.290, Z = 471.823 }

Console.WriteLine( point1 == point2); // Output: True
Console.WriteLine(point1.Equals(point2)); // Output: True

public record Point(string Name, double X, double Y, double Z);
\end{lstlisting}

          point对象的属性值在初始化后不能被修改，尝试修改会导致编译错误，所以上面代码中第2行代码被注释掉了。

          第9、10行代码中，point1和point2对象的属性值完全相同，因此它们被认为是相等的，输出结果为True。
          这也是 record 类型不同于 class 的一个重要特性，即它们的相等性是基于属性值的，而不是基于对象引用。

          record struct 基本等同于 struct 的相等性定义。

\end{itemize}


\subsection{变量的定义与 var 关键字}

\cs 中变量是前置定义的，需先指定类型，然后是变量名，最后是其值,
如语句 \verb|int i = 0;| 定义 i 为 int 类型的变量，并赋初始值为 0。

变量也可隐式定义，让编译器根据变量的值进行推定，
如语句\verb|var i = 0;| 中的 i 将由编译器推断出类型为 int 。

在程序设计中，为了保持代码简洁，定义变量时可以使用 var 关键字让编译器推断变量的数据类型。

值类型数据可以在声明时赋初始值，引用类型变量在定义时需用 new 生成实例对象或
将其它实例对象的引用赋值给它，或直接赋值为 null。

\subsection{变量的定义位置}
\cs 中变量的定义位置有三种：字段、局部变量和参数。
作为字段的变量在类或结构中定义，作为局部变量的变量在方法、构造函数、属性或索引器中定义，
\cs 中的参数是方法、构造函数、属性或索引器的参数列表中的变量。

\cs 中不存在全局变量，所有变量都必须在类或结构中定义。



\section{表达式与语句}

\subsection{表达式}
表达式由操作数 (operand) 和运算符 (operator) 构成。表达式的运算符指示对操作数适用什么样的运算。
运算符的示例包括+、-、*、/ 和 new。操作数的示例包括文本、字段、局部变量和表达式。

\begin{tabular}{l|l}
    \hline
    类别              & 说明                               \\
    \hline
    x.m             & 成员访问                             \\
    x(...)          & 方法和委托调用                          \\
    x[...]          & 数组和索引器访问                         \\
    x++             & 后增量                              \\
    x--             & 后减量                              \\
    new T(...)      & 对象和委托创建                          \\
    new T(...){...} & 使用初始值设定项创建对象                     \\
    new {...}       & 匿名对象初始值设定项                       \\
    new T[...]      & 数组创建                             \\
    typeof(T)       & 获取 T 的 System.Type 对象            \\
    checked(x)      & 在 checked 上下文中计算表达式              \\
    unchecked(x)    & 在 unchecked 上下文中计算表达式            \\
    default(T)      & 获取类型 T 的默认值                      \\
    delegate {...}  & 匿名函数（匿名方法）                       \\
    (T)x            & 将 x 显式转换为类型 T                    \\
    await x         & 异步等待 x 完成                        \\
    x is T          & 如果 x 为 T，则返回 true，否则返回 false     \\
    x as T          & 返回转换为类型 T 的 x，如果 x 不是 T 则返回 null \\
    (T x) => y      & 匿名函数（lambda 表达式）                 \\
    \hline
\end{tabular}

其它的与 C 或 \cpp 语言基本相同。


\subsection{ 语句 }
程序的操作是使用语句 (statement) 来表示与实现的。

\cs 中常见的语句有：
\begin{itemize}
    \item 声明语句 (declaration statement)

          用于声明局部变量和常量。

    \item 表达式语句 (expression statement)

          用于对表达式求值。

          可用作语句的表达式包括方法调用、使用 new 运算符的对象分配、
          使用 = 和复合赋值运算符的赋值、使用 ++ 和 -- 运算符的增量和减量运算以及 await 表达式。

    \item 选择语句 (selection statement)

          用于根据表达式的值从若干个给定的语句中选择一个来执行。
          这一组中有 if else、 switch case、 ? : 、 ??语句。

          ??表示如果左值为null，则用右值进行计算，如下代码所示：
          \begin{lstlisting}
int? a = null;
int b = a ?? -1;
Console.WriteLine(b);  // output: -1
\end{lstlisting}

    \item 迭代语句 (iteration statement)

          用于重复执行嵌入语句，常用语句有 while、do、for 和 foreach 语句。

          \cs 中常用foreach语句进行循环，示例代码如下:

          \begin{lstlisting}
static void Main(string[] args)
{
    foreach (string s in args)
    {
        Console.WriteLine(s);
    }
}
\end{lstlisting}


    \item 跳转语句 (jump statement)

          用于转移控制，常用语句有 break、continue、goto、throw、return 和 yield 语句。

    \item 异常捕捉

          try...catch 语句用于捕获在块的执行期间发生的异常，try...finally 语句用于指定终止代码，不管是否发生异常，该代码都始终要执行。

          \begin{lstlisting}
static double Divide(double x, double y)
{
    if (y == 0) throw new DivideByZeroException();
    return x / y;
}

static void Main(string[] args)
{
    try
    {
        if (args.Length != 2)
        {
            throw new Exception("Two numbers required");
        }
        double x = double.Parse(args[0]);
        double y = double.Parse(args[1]);
        Console.WriteLine(Divide(x, y));
    }
    catch (Exception e)
    {
        Console.WriteLine(e.Message);
    }
    finally
    {
        Console.WriteLine(“Good bye!”);
    }
}
\end{lstlisting}

    \item checked 语句和 unchecked 语句

          用于控制整型算术运算和转换的溢出检查上下文。

          \begin{lstlisting}
static void Main()
{
    int i = int.MaxValue;
    checked 
    {
        Console.WriteLine(i + 1);       // Exception
    }
    unchecked
    {
        Console.WriteLine(i + 1);       // Overflow
    }
}
\end{lstlisting}

    \item lock 语句

          用于获取某个给定对象的互斥锁，执行一个语句，然后释放该锁。

          \begin{lstlisting}
class Account
{
    decimal balance;
    public void Withdraw(decimal amount)
    {
        lock (this)
        {
            if (amount > balance)
            {
                throw new Exception("Insufficient funds");
            }
            balance -= amount;
        }
    }
}
\end{lstlisting}

    \item using 语句

          用于获得一个资源，执行一个语句，然后释放该资源。

          \begin{lstlisting}
static void Main()
{
    using (TextWriter w = File.CreateText("test.txt"))
    {
        w.WriteLine("Line one");
        w.WriteLine("Line two");
        w.WriteLine("Line three");
    }
}
\end{lstlisting}


\end{itemize}


\section{ \cs 特殊语法}

\subsection{泛型}
在定义类时，在类名后用尖括号括起来的类型参数列表，称之为泛型。如：
\begin{lstlisting}
public class Pair<T1, T2>
{
    public T1 first;
    public T2 second;
}
\end{lstlisting}
代码中的T1与T2为类型参数。

泛型是现代编程语言非常重要的语法。

\subsection{集合}
数组Array、列表List、词典Dictionary、栈Stack、队列Queue
是程序设计中常用的数据结构，在 \cs 中均以类库形式提供，
在\cs 程序设计中可以直接使用。

在使用这些集合类时，应尽量使用泛型集合,如动态数组：List<int>

\subsection{委托}
委托类似于C、\cpp 中的函数指针。语法为：

访问修饰符 delegate 返回类型 委托名(参数序列);

从语法上可以看出，delegate之后实际上就是函数的定义。

.Net中预定义了Func和Action委托。

Func委托代表有返回值的委托，参数最多有16个，最后一个参数必须是返回参数TResult，语法形式为：

public delegate TResult Func<in T1, ..., in T15, out TResult>(T1 arg1, ..., T15 arg15);

Action委托代表无有返回值的委托，参数最多有16个，语法形式为：

public delegate void Action<in T1,...,in T16>(T1 arg1, ..., T16 arg2);

\subsection{Lambda语句}
Lambda表达式主要用于简化委托的编写形式，是一种可用于创建委托或表达式树类型的匿名函数。语法为：

(输入参数列表) => \{表达式或语句块\}

Lambda表达式实质是一个匿名函数，如下代码所示：
\begin{lstlisting}
List<int> list = new List<int>(){1,2,3,4,5,6,7,8,10,11,12,13,14,15};
var q = list.Were( x => x%2 == 0);
\end{lstlisting}
list.Where函数中调用的就是Lambda表达式。

\subsection{迭代器}

迭代器用于逐步迭代集合(如列表和数组)。

迭代器方法或 get 访问器可对集合执行自定义迭代。 迭代器方法使用 yield return 语句返回元素，每次返回一个。 到达 yield return 语句时，会记住当前在代码中的位置。 下次调用迭代器函数时，将从该位置重新开始执行。

通过 foreach 语句或 LINQ 查询从客户端代码中使用迭代器。

如下所示代码，将输出 3 5 8

\begin{lstlisting}
using System;

namespace ConsoleApp1
{
    class Program
    {
        static void Main(string[] args)
        {
            foreach (int number in SomeNumbers())
            {
                Console.Write(number.ToString() + " ");
            }
            // Output: 3 5 8
            Console.ReadKey();

        }

        public static System.Collections.IEnumerable SomeNumbers()
        {
            yield return 3;
            yield return 5;
            yield return 8;
        }
    }
}
\end{lstlisting}

迭代器方法或 get 访问器的返回类型可以是 IEnumerable、IEnumerable<T>、IEnumerator 或 IEnumerator<T>。

可以使用 yield break 语句来终止迭代。

yield语句

\begin{lstlisting}
static IEnumerable<int> Range(int from, int to) {
    for (int i = from; i < to; i++) {
        yield return i;
    }
    yield break;
}
static void Main()
{
    foreach (int x in Range(-10,10)) {
        Console.WriteLine(x);
    }
}
\end{lstlisting}



\section{小结}
这一章我们学习了 \cs 基本语法：数据与语句，数据分为数值类型与引用类型，
数据定义时可以使用var关键字让编译器推断变量的数据类型，还可以使用泛型。

\cs 语句相对于C、\cpp 而言，出现了一些新的如lambda语句、迭代器等，
这些是许多现代语言中必不可少的语言特征，
我们将在后续课程中将其继续学习与应用。

\section{上机实验}
\begin{enumerate}
    \item 冒泡排序 (Bubble Sort)

          冒泡排序是一种较简单的排序算法。它重复地走访过要排序的元素列，
          依次比较两个相邻的元素，如果顺序错误就把他们交换过来。
          走访元素的工作是重复地进行直到没有相邻元素需要交换，也就是说该元素列已经排序完成。
          这个算法的名字由来是因为越小的元素会经由交换慢慢“浮”到数列的顶端（升序或降序排列），
          就如同碳酸饮料中二氧化碳的气泡最终会上浮到顶端一样，故名“冒泡排序”。

          原理如下：

          比较相邻的元素。如果第一个比第二个大，就交换他们两个。

          针对所有的元素重复以上的步骤，除了最后一个。

          持续每次对越来越少的元素重复上面的步骤，直到没有任何一对数字需要比较；

    \item 快速排序 (Quick Sort)

          基本思想：通过一趟排序将待排记录分隔成独立的两部分，其中一部分记录的关键字均比另一部分的关键字小，则可分别对这两部分记录继续进行排序，以达到整个序列有序。

          算法描述：

          快速排序使用分治法来把一个串（list）分为两个子串（sub-lists）。具体算法描述如下：

          从数列中挑出一个元素，称为 “基准”（pivot）；

          重新排序数列，所有元素比基准值小的摆放在基准前面，所有元素比基准值大的摆在基准的后面（相同的数可以到任一边）。在这个分区退出之后，该基准就处于数列的中间位置。这个称为分区（partition）操作；

          递归地（recursive）把小于基准值元素的子数列和大于基准值元素的子数列排序。

\end{enumerate}

